\documentclass[12pt]{article}

% Layout.
\usepackage[top=1.2in, bottom=0.75in, left=1in, right=1in, headheight=1.0in, headsep=0pt]{geometry}

% Fonts.
\usepackage{mathptmx}
\usepackage[scaled=0.86]{helvet}
\renewcommand{\emph}[1]{\textsf{\textbf{#1}}}

% TiKZ.
\usepackage{tikz, pgfplots}
\usetikzlibrary{calc}
\pgfplotsset{my style/.append style={axis x line=middle, axis y line=middle, xlabel={$x$}, ylabel={$y$}}}
\pgfplotsset{compat=1.16}

% Misc packages.
\usepackage{amsmath,amssymb,latexsym}
\usepackage{graphicx}
\usepackage{array}
\usepackage{xcolor}
\usepackage{multicol}

% Commands to set various header/footer components.
\makeatletter
\def\doctitle#1{\gdef\@doctitle{#1}}
\doctitle{Use {\tt\textbackslash doctitle\{MY LABEL\}}.}
\def\docdate#1{\gdef\@docdate{#1}}
\docdate{Use {\tt\textbackslash docdate\{MY DATE\}}.}
\def\doccourse#1{\gdef\@doccourse{#1}}
\let\@doccourse\@empty
\def\docscoring#1{\gdef\@docscoring{#1}}
\let\@docscoring\@empty
\def\docversion#1{\gdef\@docversion{#1}}
\let\@docversion\@empty
\makeatother

% Headers and footers layout.
\makeatletter
\usepackage{fancyhdr}
\pagestyle{fancy}
\fancyhf{} % Clears all headers/footers.
\lhead{\emph{\@doctitle\hfill\@docdate} \medskip
\ifnum \value{page} > 1\relax\else\\
\emph{Name: \rule{3.5in}{1pt}\hfill \@docscoring}
\fi}

\rfoot{\emph{\@docversion}}
\lfoot{\emph{\@doccourse}}
\cfoot{\emph{\thepage}}
\renewcommand{\headrulewidth}{0pt}%
\makeatother

% Paragraph spacing
\parindent 0pt
\parskip 6pt plus 1pt

% A problem is a section-like command. Use \problem{5} for a problem worth 5 points.
\newcounter{probcount}
\newcounter{subprobcount}
\setcounter{probcount}{0}

\newcommand{\problem}[1]{%
\par
\addvspace{4pt}%
\setcounter{subprobcount}{0}%
\stepcounter{probcount}%
\makebox[0pt][r]{\emph{\arabic{probcount}.}\hskip1ex}\emph{[#1 points]}\hskip1ex}

\newcommand{\thesubproblem}{\emph{\alph{subprobcount}.}}

% like \problem but with name
\newcommand{\nameprob}[2]{%
\par
\addvspace{4pt}%
\setcounter{subprobcount}{0}%
\stepcounter{probcount}%
\makebox[0pt][r]{\emph{#1.}\hskip1ex}\emph{[#2 points]}\hskip1ex}

% Subproblems are an enumerate-like environment with a consistent
% numbering scheme.  Use \begin{subproblems}\item...\item...\end{subproblems}
\newenvironment{subproblems}{%
\begin{enumerate}%
\setcounter{enumi}{\value{subprobcount}}%
\renewcommand{\theenumi}{\emph{\alph{enumi}}}}%
{\setcounter{subprobcount}{\value{enumi}}\end{enumerate}}

% Blanks for answers in normal and math mode.
\newcommand{\blank}[1]{\rule{#1}{0.75pt}}
\newcommand{\mblank}[1]{\underline{\hspace{#1}}}
\def\emptybox(#1,#2){\framebox{\parbox[c][#2]{#1}{\rule{0pt}{0pt}}}}

% Misc.
\renewcommand{\d}{\displaystyle}
\newcommand{\ds}{\displaystyle}


\doctitle{Math 252: Quiz 10}
\docdate{23 Apr 2026}
\doccourse{}
\docversion{}
\docscoring{\fbox{{\LARGE \strut}\blank{0.8in} / 25}}

\begin{document}
30 minutes.  No aids (book, notes, calculator, internet, etc.) are permitted.  Show all work and use proper notation for full credit.  Put answers in reasonably-simplified form.  25 points possible.

\problem{8} Write the MacLaurin Series of each function. 

	\begin{subproblems}
    %HW 6.4 #175
	\item $f(x) = (1+x^2)^{1/3}$
	\vfill
    %HW 6.4 #231
	\item $\displaystyle \cosh(x) = \frac{e^x+e^{-x}}{2}$. Express your answer as {\bf one series}.
	\vfill
	\end{subproblems}

\newpage 

\problem{8}
Consider the parametric equations $x=2t^2$,\quad  $y=t^4+1$,\quad  $t\in[0,2]$.
	\begin{subproblems}
    %HW 7.1 #5
	\item  Eliminate the parameter and sketch the graph.
	\vfill \vfill 
			
	\item Find $\ds \frac{d^2y}{dx^2}$ using the result of part {\bf a}.
	\vfill

    \item Find $\ds \frac{d^2y}{dx^2}$ using the parametric formula, and check that it gives the same result as in part {\bf b}.
    \vfill \vfill
	
\end{subproblems}	
	
	%\vspace{.5in}

\newpage 


\problem{9} 
Consider the parametric equations $x=t-\sin t$,\quad $y=1-\cos t$, \quad $t\in [0,2\pi]$.
	\begin{subproblems}
    \item Find the equation of the tangent line at $t=\frac{\pi}{2}$.
	\vfill\vfill
    
	\item  Set up (but do not evaluate) an integral for the area under the parametric curve.
	\vfill
			
	\item Set up (but do not evaluate) an integral for the arc length of the parametric curve.
	\vfill
	
\end{subproblems}	

\newpage 
\textbf{\textsf{Extra Credit. [1 point]}} {\large\strut} \, Evaluate the integral in \textbf{Problem 3 c.}
	\vfill \vfill \vfill
    
\noindent \hrule
\medskip
\centerline{\footnotesize \textsc{blank space}}
\vfill \vfill
\end{document}